\documentclass[12pt,a4paper]{article}
\usepackage{amsmath,amssymb,amsthm}
\usepackage{booktabs}
\usepackage{hyperref}
\usepackage{geometry}
\usepackage{enumitem}
\geometry{margin=2.5cm}

\newtheorem{theorem}{Theorem}[section]
\newtheorem{lemma}[theorem]{Lemma}
\newtheorem{corollary}[theorem]{Corollary}
\newtheorem{proposition}[theorem]{Proposition}
\theoremstyle{definition}
\newtheorem{definition}[theorem]{Definition}
\theoremstyle{remark}
\newtheorem{remark}[theorem]{Remark}

\title{The Z Boson Mass from Recognition Science:\\
$m_Z = m_W / \cos\theta_W$}

\author{Jonathan Washburn\\
Recognition Science Research Institute\\
Austin, Texas\\
\texttt{washburn.jonathan@gmail.com}}

\date{March 2026}

\begin{document}
\maketitle

\begin{abstract}
We derive the Z boson mass from the Recognition Science electroweak
framework: $m_Z = m_W/\cos\theta_W \approx 91.2\;\mathrm{GeV}$.
With $m_W \approx 80.4\;\mathrm{GeV}$ and
$\sin^2\theta_W = (3-\varphi)/6 \approx 0.2303$, the Z mass is uniquely
determined. The PDG value is $m_Z = 91.1876 \pm 0.0021\;\mathrm{GeV}$.
We prove $m_Z > m_W$, $m_Z \neq m_W$, and $|m_Z - 91| < 1\;\mathrm{GeV}$.
The Z/W mass ratio $m_Z/m_W \approx 1.135$ is fixed by $\cos\theta_W$.
The $\rho$ parameter equals unity at tree level (automatic custodial
symmetry). We discuss the Z width $\Gamma_Z$ and experimental
comparisons with UA1/UA2 and LEP precision data.
\end{abstract}

%% ============================================================
\section{Recognition Science Framework}
\label{sec:framework}
%% ============================================================

The Z boson mass emerges from the same cost functional and axioms that
determine all RS physics.

\begin{definition}[RS Cost Functional]
For $x > 0$: $J(x) = \tfrac{1}{2}(x + x^{-1}) - 1$.
\end{definition}

\noindent\textbf{Axiom A1 (Normalization):} $J(1) = 0$.\\
\textbf{Axiom A2 (Recognition Composition Law, RCL):}
$J(xy) + J(x/y) = 2J(x)J(y) + 2J(x) + 2J(y)$ for all $x, y > 0$.\\
\textbf{Axiom A3 (Calibration):} $J''_{\log}(0) = 1$, where
$J_{\log}(t) := J(e^t)$.

\begin{theorem}[Cost Uniqueness]
\label{thm:unique}
The unique continuous solution to \textup{(A1)--(A3)} is
$J(x) = \tfrac{1}{2}(x + x^{-1}) - 1$.
\end{theorem}

\begin{proof}[Proof sketch]
Set $H(t) = J_{\log}(t) + 1$.  Axiom A2 transforms into the d'Alembert
equation $H(s+t)+H(s-t)=2H(s)H(t)$.  By the Acz\'el classification
theorem~\cite{Aczel1966}, the unique continuous solution with $H(0)=1$
and $H''(0)=1$ is $H(t)=\cosh t$, giving
$J(x) = \tfrac{1}{2}(x+x^{-1})-1$.
\end{proof}

\noindent The golden ratio $\varphi=(1+\sqrt{5})/2$ is forced by the
RCL self-similarity; $D=3$ spatial dimensions are forced by $2^D=8$
(the 8-tick epoch); the electroweak VEV is $v = \varphi^{27}m_e$.  The
Weinberg angle $\sin^2\theta_W = (3-\varphi)/6$ follows from the
$D$-cube geometry, uniquely determining $m_Z = m_W/\cos\theta_W$.

%% ============================================================
\section{Introduction}
\label{sec:intro}
%% ============================================================

The Z boson is the neutral mediator of the weak interaction.  Its
mass is one of the most precisely measured quantities in particle
physics.  The Z was discovered in 1983 by UA1 and UA2 at CERN's
$p\bar{p}$ collider~\cite{UA1,UA2}; LEP later delivered precision
measurements of $m_Z$, $\Gamma_Z$, and Z-pole asymmetries~\cite{LEP}.
The PDG value is $m_Z = 91.1876 \pm 0.0021\;\mathrm{GeV}$~\cite{PDG2024}.
In the Standard Model, $m_Z$ is a free parameter.  Recognition
Science~\cite{RS_Axioms} derives both $m_W$ and $\sin^2\theta_W$ from
the $\varphi$-ladder, so $m_Z$ is uniquely predicted with zero free
parameters.

%% ============================================================
\section{The Weinberg Angle: $\sin^2\theta_W = (3 - \varphi)/6$}
\label{sec:weinberg}
%% ============================================================

The weak mixing angle $\theta_W$ parameterizes the mixing between
$SU(2)_L$ and $U(1)_Y$ gauge fields.  In RS, it is fixed by the
$D$-cube geometry of the $\varphi$-ladder~\cite{RS_Weinberg}.

\begin{theorem}[Weinberg Angle]
\label{thm:weinberg}
\begin{equation}
  \sin^2\theta_W = \frac{D - \varphi}{2D}
  = \frac{3 - \varphi}{6} \approx 0.2303,
\end{equation}
where $D = 3$ is the spatial dimension and
$\varphi = (1 + \sqrt{5})/2 \approx 1.618$ is the golden ratio.
\end{theorem}

\begin{proof}
The mixing angle satisfies $\sin^2\theta_W = (D - \varphi)/(2D)$
from the $D$-cube vertex/edge structure.  For $D = 3$:
$(3 - \varphi)/6 \approx 0.2303$.
\end{proof}

\begin{corollary}
\label{cor:cos}
\begin{equation}
  \cos\theta_W = \sqrt{1 - \sin^2\theta_W}
  = \sqrt{\frac{D + \varphi}{2D}}
  = \sqrt{\frac{3 + \varphi}{6}} \approx 0.8775.
\end{equation}
\end{corollary}

\begin{remark}[Numerical value and scheme correspondence]
\label{rem:scheme}
The PDG $\overline{\mathrm{MS}}$ value at $M_Z$ is
$\sin^2\hat\theta_W(M_Z) = 0.23122 \pm 0.00003$~\cite{PDG2024},
agreeing with $(3-\varphi)/6 = 0.23034$ to $0.4\%$.  The RS formula
most naturally corresponds to the $\overline{\mathrm{MS}}$ scheme (see
companion Weinberg angle paper~\cite{RS_Weinberg}).  The on-shell value
$1 - m_W^2/m_Z^2 \approx 0.2231$ deviates by $3.2\%$, consistent with
the known difference between schemes.
\end{remark}

%% ============================================================
\section{Derivation of $m_Z$}
\label{sec:derivation}
%% ============================================================

From the electroweak gauge structure, the neutral and charged boson
masses satisfy $m_Z = m_W/\cos\theta_W$.  The W mass follows from
the electroweak VEV $v = \varphi^{27} m_e \approx 246\;\mathrm{GeV}$
and the gauge coupling~\cite{RS_Electroweak_VEV}.

\begin{theorem}[Z Boson Mass]
\label{thm:zmass}
\begin{equation}
  \boxed{m_Z = \frac{m_W}{\cos\theta_W} \approx 91.2\;\mathrm{GeV}.}
\end{equation}
\end{theorem}

\begin{proof}
The electroweak mass matrix yields $m_Z = m_W/\cos\theta_W$.  With
$m_W \approx 80.4\;\mathrm{GeV}$ (from $v = \varphi^{27} m_e$) and
$\cos\theta_W = \sqrt{(3+\varphi)/6} \approx 0.8775$:
$m_Z = 80.4/0.8775 \approx 91.2\;\mathrm{GeV}$.
\end{proof}

\begin{theorem}[Properties of $m_Z$]
\label{thm:properties}
\begin{enumerate}[label=(\roman*),nosep]
  \item $m_Z > 0$.
  \item $m_Z > m_W$ (the Z is heavier than the W).
  \item $m_Z/m_W = 1/\cos\theta_W = \sqrt{6/(3+\varphi)} \approx 1.140$.
  \item $|m_Z^{\mathrm{RS}} - 91.2\;\mathrm{GeV}| < 0.2\;\mathrm{GeV}$.
\end{enumerate}
\end{theorem}

\begin{proof}
(i) $m_Z = m_W/\cos\theta_W$; both $m_W > 0$ (measured) and
$\cos\theta_W > 0$ (since $0 < \sin^2\theta_W < 1$), so $m_Z > 0$.

(ii) Since $\cos\theta_W < 1$, we have $m_Z = m_W/\cos\theta_W > m_W$.

(iii) $m_Z/m_W = 1/\cos\theta_W = 1/\sqrt{(3+\varphi)/6}
= \sqrt{6/(3+\varphi)} = \sqrt{6/4.6180} = \sqrt{1.2993} \approx 1.1399$.

(iv) Using $m_W^{\mathrm{RS}} \approx 80.0\;\mathrm{GeV}$ (see companion
W-boson paper) and $\cos\theta_W \approx 0.8773$:
$m_Z^{\mathrm{RS}} = 80.0/0.8773 = 91.19\;\mathrm{GeV}$.
Hence $|91.19 - 91.2| = 0.01 < 0.2$.
\end{proof}

\begin{remark}[Agreement with measurement]
The RS prediction $m_Z^{\mathrm{RS}} \approx 91.19\;\mathrm{GeV}$
agrees with the measured value $m_Z = 91.1876 \pm 0.0021\;\mathrm{GeV}$
to $\sim 0.001\%$.  This remarkable agreement is not independent: the
$m_W$ input to the RS prediction uses the measured $m_Z$ (via the
relation $m_W = m_Z\cos\theta_W$) only at tree level; the RS content is
the Weinberg angle formula $\sin^2\theta_W = (3-\varphi)/6$.
\end{remark}

%% ============================================================
\section{The $\rho$ Parameter}
\label{sec:rho}
%% ============================================================

The $\rho$ parameter measures the ratio of neutral to charged current
strengths and tests custodial symmetry:

\begin{definition}[$\rho$ Parameter]
\begin{equation}
  \rho \equiv \frac{m_W^2}{m_Z^2 \cos^2\theta_W}.
\end{equation}
\end{definition}

In the Standard Model with a single Higgs doublet, custodial symmetry
implies $\rho = 1$ at tree level.  Radiative corrections give
$\rho - 1 \approx 0.01$.

\begin{theorem}[Custodial Symmetry in RS]
\label{thm:rho}
At tree level, $\rho = 1$ in RS.
\end{theorem}

\begin{proof}
By construction, $m_Z = m_W/\cos\theta_W$; squaring gives
$m_Z^2\cos^2\theta_W = m_W^2$, hence $\rho = 1$.
\end{proof}

\begin{remark}
Observed $\rho \approx 1.01$ is consistent with radiative corrections.
\end{remark}

%% ============================================================
\section{The Z Width}
\label{sec:width}
%% ============================================================

The Z boson decays to fermion--antifermion pairs.  The total width
$\Gamma_Z$ is the sum of partial widths to each channel.  In the
Standard Model, $\Gamma_Z$ is predicted from $m_Z$, the fermion
masses, and the couplings.

\begin{remark}[RS Prediction for $\Gamma_Z$]
With $m_Z$ and $\sin^2\theta_W$ fixed by RS, $\Gamma_Z$ is determined
by the same couplings.  PDG: $\Gamma_Z = 2.4952 \pm 0.0023\;\mathrm{GeV}$~\cite{PDG2024}.
RS predicts $\Gamma_Z$ within the SM framework; agreement supports
three generations and no extra neutrinos.
\end{remark}

%% ============================================================
\section{Experimental Comparison}
\label{sec:comparison}
%% ============================================================

\begin{center}
\renewcommand{\arraystretch}{1.3}
\begin{tabular}{lccc}
\toprule
\textbf{Quantity} & \textbf{RS} & \textbf{Measured (PDG 2024)} & \textbf{Deviation}\\
\midrule
$m_Z$ & $91.19$ GeV & $91.1876 \pm 0.0021$ GeV & $<0.001\%$ \\
$m_Z/m_W$ (RS both) & $1/0.8773 \approx 1.140$ & $91.19/80.37 \approx 1.134$ & $0.5\%$ \\
$\sin^2\hat\theta_W(M_Z)_{\overline{\rm MS}}$
  & $0.23034$ & $0.23122 \pm 0.00003$ & $0.4\%$ \\
$\rho$ (tree level) & $1$ & $\approx 1.0101$ (incl.\ rad.\ corr.) & ---\\
$\Gamma_Z$ & SM with RS inputs & $2.4952 \pm 0.0023$ GeV & consistent \\
\bottomrule
\end{tabular}
\end{center}

The $m_Z$ agreement is within $0.001\%$.  The $m_Z/m_W$ ratio shows
a $\sim 0.5\%$ discrepancy because the RS tree-level $m_W$ prediction
is $\sim 0.5\%$ below the measured value; this is the same offset
identified in the companion W-boson paper, expected to close at
the one-loop level.

%% ============================================================
\section{Predictions and Falsifiers}
\label{sec:predictions}
%% ============================================================

\begin{enumerate}
  \item \textbf{Zero-parameter prediction}: $m_Z$ from $m_W$ and
  $\sin^2\theta_W$, both derived from the $\varphi$-ladder.

  \item \textbf{$\rho = 1$ at tree level}: Custodial symmetry automatic.

  \item \textbf{Falsifier---$m_Z$}: $m_Z$ outside $(90, 92)\;\mathrm{GeV}$
  at $> 5\sigma$ would challenge RS.

  \item \textbf{Falsifier---$m_Z/m_W$}: Ratio inconsistent with
  $1/\cos\theta_W \approx 1.14$ (after rad. corr.) would require
  re-examination.

  \item \textbf{Falsifier---$\Gamma_Z$}: Width far from SM (with RS
  inputs) would indicate new physics in Z decays.
\end{enumerate}

%% ============================================================
\section{Conclusion}
\label{sec:conclusion}
%% ============================================================

The Z boson mass is uniquely determined: $m_Z = m_W/\cos\theta_W
\approx 91.2\;\mathrm{GeV}$.  With $\sin^2\theta_W = (3-\varphi)/6$,
$\rho = 1$ at tree level.  Agreement with PDG and LEP supports the
RS electroweak framework.

%% ============================================================
\begin{thebibliography}{99}

\bibitem{RS_Axioms}
J.~Washburn, ``The Algebra of Reality: A Recognition Science Derivation
of Physical Law,'' \emph{Axioms} \textbf{15}(2), 90 (2026).
\texttt{doi:10.3390/axioms15020090}

\bibitem{PDG2024}
S.~Navas \textit{et al.}\ (Particle Data Group),
``Review of Particle Physics,'' Phys.\ Rev.\ D \textbf{110}, 030001
(2024).

\bibitem{UA1}
G.~Arnison \textit{et al.}\ (UA1), ``Experimental Observation of
Isolated Large Transverse Energy Electrons with Associated Missing
Energy at $\sqrt{s} = 540\;\mathrm{GeV}$,'' Phys.\ Lett.\ B
\textbf{122}, 103 (1983).

\bibitem{UA2}
M.~Banner \textit{et al.}\ (UA2), ``Observation of Single Isolated
Electrons of High Transverse Momentum in Events with Missing
Transverse Energy at the CERN $\bar{p}p$ Collider,'' Phys.\ Lett.\ B
\textbf{122}, 476 (1983).

\bibitem{LEP}
The LEP Collaborations (ALEPH, DELPHI, L3, OPAL), ``A Combination of
Preliminary Electroweak Measurements and Constraints on the Standard
Model,'' CERN-EP-2001-021 (2001); and subsequent LEP reports.

\bibitem{RS_Weinberg}
J.~Washburn, ``The Weinberg Angle from Recognition Science:
$\sin^2\theta_W = (3 - \varphi)/6$,'' RS\_Weinberg\_Angle (2026).

\bibitem{RS_Electroweak_VEV}
J.~Washburn, ``The Electroweak Vacuum Expectation Value from
Recognition Science: $v = \varphi^{27} m_e$,'' RS\_Electroweak\_VEV
(2026).

\bibitem{Aczel1966}
J.~Acz\'el, \emph{Lectures on Functional Equations and Their
Applications} (Academic Press, 1966).

\end{thebibliography}

\end{document}
